\ticket{19}{Обучение без учителя и кластеризация текстов}

\begin{examplan}
    \item Дать определение обучения без учителя.
    \item Описать задачу кластеризации и метрики качества.
    \item Объяснить метод k-средних.
    \item Объяснить иерархическую кластеризацию.
    \item Рассказать о кластеризации текстов и признаках.
\end{examplan}

\subsection*{Короткая версия для письменного ответа}

\textbf{Обучение без учителя} работает с неразмеченными данными и ищет скрытую структуру: кластеры, факторы, аномалии, ассоциации или представления меньшей размерности.

\textbf{Кластеризация} --- разбиение объектов на группы так, чтобы объекты внутри кластера были похожи, а объекты из разных кластеров различались.

\subsection*{Оценка качества кластеризации}

Оценка сложна, потому что правильных меток обычно нет.

\textbf{Внутренние меры} используют только данные:
\begin{itemize}
    \item silhouette score: близко к 1 --- объект хорошо лежит в своем кластере, около 0 --- на границе, меньше 0 --- вероятно не в своем кластере;
    \item индекс Данна: отношение межкластерного расстояния к внутрикластерному;
    \item индекс Дэвиса--Болдина: меньше --- лучше.
\end{itemize}

\textbf{Внешние меры} используют известные метки, если они есть: Rand Index, Adjusted Rand Index, purity, F-мера.

\subsection*{Метод k-средних}

\textbf{k-means} разбивает объекты на $k$ кластеров, минимизируя сумму квадратов расстояний до центроидов.

Алгоритм:
\begin{enumerate}
    \item выбрать $k$ начальных центроидов;
    \item отнести каждый объект к ближайшему центроиду;
    \item пересчитать центроид как среднее объектов кластера;
    \item повторять до стабилизации или лимита итераций.
\end{enumerate}

Особенности:
\begin{itemize}
    \item число кластеров $k$ задается заранее;
    \item результат зависит от начальной инициализации;
    \item алгоритм сходится к локальному минимуму;
    \item хорошо работает для компактных примерно сферических кластеров;
    \item плохо работает с кластерами сложной формы, разной плотности и выбросами.
\end{itemize}

Выбор $k$: метод локтя, silhouette score, экспертная интерпретация.

\subsection*{Иерархическая кластеризация}

\textbf{Иерархическая кластеризация} строит дерево кластеров --- дендрограмму.

Виды:
\begin{itemize}
    \item \textbf{агломеративная}: начинается с отдельных объектов и последовательно объединяет ближайшие кластеры;
    \item \textbf{дивизивная}: начинается с одного кластера и последовательно делит его.
\end{itemize}

Критерии расстояния между кластерами:
\begin{itemize}
    \item single linkage: расстояние между ближайшими объектами, может давать эффект цепочки;
    \item complete linkage: расстояние между самыми дальними объектами, дает более компактные кластеры;
    \item average linkage: среднее расстояние между объектами;
    \item метод Уорда: минимизирует рост внутрикластерной суммы квадратов.
\end{itemize}

Плюс: не нужно заранее фиксировать число кластеров. Минус: высокая вычислительная сложность и необратимость ранних объединений.

\subsection*{Кластеризация текстов}

Перед кластеризацией тексты преобразуют в векторы.

Этапы:
\begin{enumerate}
    \item токенизация;
    \item приведение к нижнему регистру;
    \item удаление стоп-слов;
    \item лемматизация или стемминг;
    \item построение признаков: BoW, TF-IDF, n-граммы, эмбеддинги документов.
\end{enumerate}

Для разреженных TF-IDF-векторов часто используют \textbf{косинусную близость}, потому что важнее направление вектора, а не его длина. Для плотных эмбеддингов могут использоваться косинусная или евклидова метрика.

\begin{important}
    \item В обучении без учителя нет заранее заданных правильных ответов.
    \item Кластеризация ищет группы похожих объектов.
    \item k-means быстрый, но требует выбрать $k$ и чувствителен к инициализации.
    \item Иерархическая кластеризация дает дендрограмму и позволяет выбрать число кластеров после построения.
    \item Для текстов особенно важны векторизация и метрика близости.
\end{important}

